\documentclass[12pt]{article}

\usepackage[spanish]{babel}
\usepackage[utf8]{inputenc}
\usepackage[T1]{fontenc}
\usepackage{geometry}
\usepackage{array}
\usepackage{longtable}
\usepackage{hyperref}

\geometry{letterpaper,margin=2.5cm}

% Comandos para líneas centradas a lo ancho (más cortas para evitar overfull)
\newcommand{\CampoCentrado}[2]{%
	\noindent\makebox[\textwidth][c]{#1\ \underline{\makebox[0.7\textwidth][c]{#2}}}%
}

\newcommand{\CampoCentradoVacio}[1]{%
	\noindent\makebox[\textwidth][c]{#1\ \underline{\makebox[0.7\textwidth][c]{\hspace{1pt}}}}%
}

\begin{document}
	
	% Encabezado institucional
	\begin{center}
		{\bfseries UNIVERSIDAD AUTÓNOMA METROPOLITANA}\\[4pt]
		{\bfseries UNIDAD CUAJIMALPA}\\[4pt]
		DIVISIÓN DE CIENCIAS DE LA COMUNICACIÓN Y DISEÑO\\[4pt]
		{\bfseries LICENCIATURA EN TECNOLOGÍAS DE LA INFORMACIÓN}
	\end{center}
	
	\vspace{1em}
	
	\begin{table}[h]
		\renewcommand{\arraystretch}{1.2}
		\begin{tabular}{@{}l@{\hspace{0.5cm}}l@{}}
			PROFESOR: &
			\underline{\makebox[0.7\textwidth][c]{\MakeUppercase{Juan Manuel Alvarado López}}} \\[0.3em]
			ADSCRITO AL: &
			{\small \underline{\makebox[0.7\textwidth][l]{\textbf{DEPARTAMENTO DE TECNOLOGÍAS DE LA INFORMACIÓN}}}} \\[0.3em]
			OTRO: &
			\underline{\makebox[0.7\textwidth][c]{ }} \\[0.3em]
			UEA: &
			\underline{\makebox[0.7\textwidth][c]{\MakeUppercase{Integración de sistemas}}} \\[0.3em]
		\end{tabular}
	\end{table}
	
	
	% Objetivo
	\noindent{OBJETIVO:} Conocer los principales problemas que surgen en la integración de distintos sistemas de información, destacando los aspectos referentes a la integración de sus datos, su integración funcional e interoperatividad, así como la gestión del conocimiento para su integración semántica.
	
	\vspace{1em}
	
	% HORARIO en tabla
	\section*{Horario}
	
	\begin{center}
		\renewcommand{\arraystretch}{1.3}
		\begin{tabular}{|c|c|c|}
			\hline
			\textbf{Día} & \textbf{Horario} & \textbf{Aula} \\
			\hline
			\underline{\makebox[2.8cm][c]{Martes}} &
			\underline{\makebox[2.8cm][c]{14--16}} &
			\underline{\makebox[2.8cm][c]{L-522}} \\
			\hline
			\underline{\makebox[2.8cm][c]{Jueves}} &
			\underline{\makebox[2.8cm][c]{13--16}} &
			\underline{\makebox[2.8cm][c]{L-522}} \\
			\hline
%			\underline{\makebox[2.8cm][c]{\ }} &
%			\underline{\makebox[2.8cm][c]{\ }} &
%			\underline{\makebox[2.8cm][c]{\ }} \\
%			\hline
		\end{tabular}
	\end{center}
	
	\vspace{1em}
	
	% PLAN
	\section*{Plan}
	
	\setlength{\extrarowheight}{2pt}
	
	\begin{longtable}{|>{\centering\arraybackslash}p{1.5cm}|p{6.7cm}|p{6.7cm}|}
		\hline
		\textbf{Semana} & \textbf{Tema} & \textbf{Actividades} \\
		\hline
		\endfirsthead
		\hline
		\textbf{Semana} & \textbf{Tema} & \textbf{Actividades} \\
		\hline
		\endhead
		
		1 &
		- Presentación del curso, reglas, evaluación, proyecto final 
		\newline- Panorama histórico: de integración punto a punto a SOA, servicios web SOAP/XML y surgimiento de REST.
		\newline- Concepto actual de integración: monolito vs SOA vs microservicios vs eventos.
		&
		- Mini-mapa conceptual en papel o Jamboard/Miro sobre tipos de integración (datos, funcional, presentación). (No entregable.) \\
		\hline
		
		2 &
		- Introducción formal a la integración de sistemas (tipos: datos, funcional, semántica).
		\newline- Pila clásica de servicios web SOAP: XML, WSDL, UDDI, WS-*, contratos estrictos.
		\newline- Casos actuales donde SOAP sigue siendo relevante (banca, gobierno, salud) vs APIs modernas.
		&
		- Texto de 1–1.5 cuartillas describiendo un caso real (o encontrado en la web) donde se use SOAP, indicando por qué tiene sentido allí (transacciones, regulación, legado, etc.). (Entregable) \\
		\hline
		
		3 &
		- XML para servicios web: estructura básica, esquemas XSD a nivel conceptual.
		\newline- Estructura de un mensaje SOAP: Envelope, Header, Body, Fault; comparación con mensajes REST/JSON ligeros.
		\newline- Diferencias clave SOAP vs REST (formato, estado, seguridad). &
		- Ejercicio en clase: dado un mensaje REST/JSON sencillo, escribir el “equivalente” como SOAP/XML muy básico (ignorando detalles finos de WSDL). (No entregable.) \\
		\hline
		
		4 &
		- Análisis y diseño para integración: requisitos de integración (fuentes de datos, latencia, seguridad).
		\newline- Transición conceptual: de servicios SOAP (WSDL) a diseño de APIs REST (OpenAPI) y coexistencia SOAP/REST.
		\newline- Introducción a OpenAPI/Swagger como “WSDL moderno” para REST.
		&
		- Usar Swagger online (gratuito) en clase para definir un contrato simple de API (por ejemplo, catálogo de productos: GET/POST/PUT/DELETE /products). (No entregable)
		\newline- Elegir en equipos (2–3 alumnos) el proyecto final (ej.: integración de pagos, reservas, catálogo, etc.) y entregar título + 3–4 líneas de descripción. (Entregable). \\
		\hline
		
		5 &
		- Buenas prácticas de diseño de APIs REST: recursos, URIs, verbos HTTP, manejo de errores.
		\newline- Estándares de datos: JSON y XML
		\newline- Introducción a otros estilos: gRPC y mensajería asíncrona (solo a nivel conceptual) para contrastar con SOAP/REST.
		 &
		- Ejercicio en clase: diseño de 5–6 endpoints REST para un escenario de integración (ej.: e-commerce + sistema de facturación), con rutas y métodos HTTP. (No entregable). 
		\newline- Cuadro comparativo breve SOAP vs REST vs gRPC (máx. 1 página). (Entregable) \\
		\hline
		
		6 &
		- Examen parcial (Martes, 2h) sobre: conceptos de integración, SOAP/XML, comparación SOAP–REST y diseño básico de APIs.
		 &
		Dinámica de revisión del examen en clase  \\
		\hline
		
		7 &
		- Lenguajes de descripción para servicios web: WSDL y OpenAPI.
		\newline- Documentación y descubrimiento: UDDI, portales de APIs y catálogos de servicios modernos.
		&
		Entregable: versión 1 del contrato del proyecto (revisión rápida).  \\
		\hline
		
		8 &
		- Implementación práctica de un servicio REST sencillo (Node.js/Express o Python/FastAPI) que devuelva JSON.
		\newline- Consumo de servicios SOAP desde una app moderna: patrones de “façade” REST sobre SOAP (conceptual + ejemplo).
		\newline- Seguridad básica: HTTPS, claves API sencillas.
		 &
		- Ejercicio de diseño (no implementado): en equipo bosquejar cómo envolverían un servicio SOAP legado dentro de un API REST en su proyecto. (No entregable). \\
		\hline
		
		9 &
		- Integración semántica y Web Semántica: conceptos básicos de vocabularios, ontologías ligeras, JSON-LD, etiquetas semánticas en APIs.
		\newline- Registros/catálogos de servicios modernos: portales de APIs, documentación, términos de uso.
		\newline- Patrones de integración SOAP–REST–microservicios en entornos empresariales. 
		&
		Lista breve (6–8 puntos) con los metadatos y políticas que documentarán para API del proyecto (límite de tasa, autenticación, versión, contacto, etc.). (Entregable)\\
		\hline
		
		10 &
		- Flujos de trabajo y orquestación de servicios.
		\newline- Integración síncrona vs asíncrona: llamadas REST/gRPC vs colas/eventos, webhooks.
		\newline- Ejemplos de aplicación: e-commerce, bibliotecas digitales, SIG, sistemas de pago y notificaciones, incluyendo coexistencia SOAP/REST.
		\newline- Tendencias actuales: API-first, event-driven architecture, serverless, GraphQL, APIs internas/externas/partner.
		 &
		- Actividad en clase: diagrama de flujo de un proceso (por ejemplo, compra en línea) identificando qué llamadas serían SOAP (si existiera un core legado), cuáles REST y dónde usarían eventos. (No entregable). 
		\newline- Diagrama de integración (puede ser a mano) donde se muestre: cliente, API del equipo, posible servicio externo (SOAP o REST), base de datos/cola si aplica. (Entregable) \\
		\hline
		
		11 &
		- Examen final (Martes, 2h) centrado en: comparación SOAP/REST, análisis de arquitectura, patrones de comunicación, diseño de una solución de integración. &
		- Entregable de proyecto: 1) contrato OpenAPI final; 2) breve documento (2–3 páginas) con descripción del problema, arquitectura, decisiones de integración (incluyendo si consideran coexistencia con SOAP); 3) 1–2 capturas de funcionamiento local. Presentación/demostración corta en video (5–10 min).\\
		\hline
	\end{longtable}
	
	% BIBLIOGRAFÍA
	\section*{Bibliografía}
	
	\begingroup
	\renewcommand{\refname}{}% título vacío
	\bibliographystyle{plain}     % o el estilo que uses
	\nocite{*}                  % incluye todas las entradas del .bib
	\bibliography{bibliografia}
	\endgroup
	
	
	
	% POLÍTICAS
	\section*{Políticas del curso}
	
	\begin{enumerate}
		\item \textbf{Dinámica del curso.} El curso se evalúa con tareas, exámenes parciales y un proyecto final, con ponderaciones definidas en el programa.
		\item \textbf{Faltas permitidas.} Se permiten hasta 3 faltas; dos retardos o salir sin justificación cuentan como una inasistencia. 
		\item \textbf{Eventos del profesor.} Si hay un evento que coincida con el plan de la UEA, se indicará reposición de clase o trabajo. 
		\item \textbf{Asesorías.} Se solicitan en clase o por correo institucional, acordando día, hora y lugar. 
		\item \textbf{Ponderaciones.} \begin{itemize}
			\item Tareas 40\%; 
			\item Exámenes parciales 30\%; 
			\item Proyecto final 30\%. 
		\end{itemize}
		\item \textbf{Equivalencias numéricas.} 
		\begin{itemize}
			\item NA: 0.0 – <6; 
			\item 		S: 6.0 – <7.8; 
			\item 		B: 7.8 – <8.8; 
			\item 		MB: 8.8 – 10.0.
		\end{itemize}
		\item \textbf{Deshonestidad académica.} Copiar, plagiar o falsificar implica NA como calificación final.
	\end{enumerate}
	
	
	
	\textbf{Nota}: Si el alumno incurre en cuestiones relacionadas a la deshonestidad académica tales como: copiar, plagiar o falsificar, y se le sorprende, obtendrá automáticamente como calificación final NA en esta UEA. 
	
	***Es importante que se sigan las reglas que se mencionan en la primera clase.
	
	
	
%	\section*{Firmas}
%	
%	\vspace{2cm}
%	\noindent\rule{7cm}{0.4pt}\\
%	Profesor
%	
%	\vspace{2cm}
%	\noindent\rule{7cm}{0.4pt}\\
%	Coordinador de la Licenciatura
	
\end{document}
